\subsection{BOREL SUBALGEBRAS}

PROPOSITION 7. Let \(\mathfrak{b} = \mathfrak{h} + \mathfrak{g}^{\mathrm{P}}\) be a subalgebra of \(\mathfrak{g}\) containing \(\mathfrak{h}\) . The following conditions are equivalent:

\begin{enumerate}
\def\labelenumi{(\roman{enumi})}
\item
  \(\mathfrak{b}\) is a maximal solvable subalgebra of \(\mathfrak{g}\) ;
\item
  there exists a chamber C of R such that \(P = R_{+}(C)\) ;
\item
  \(\mathrm{P} \cap (-\mathrm{P}) = \varnothing\) and \(\mathrm{P} \cup (-\mathrm{P}) = \mathrm{R}\) .
\item
  \(\Longrightarrow\) (ii): If \(\mathfrak{b}\) is solvable, \(P \cap (-P) = \varnothing\) . Then there exists a chamber \(C\) of \(R\) such that \(P \subset R_{+}(C)\) (Chap. VI, §1, no. 7, Prop. 22). Then \(\mathfrak{h} + \mathfrak{g}^{R_{+}(C)}\) is a solvable subalgebra of \(\mathfrak{g}\) containing \(\mathfrak{b}\) , hence equal to \(\mathfrak{b}\) if \(\mathfrak{b}\) is maximal.
\item
  \(\Longrightarrow\) (iii): This is clear.
\item
  \(\Longrightarrow\) (i): Assume that \(P \cap (-P) = \varnothing\) and that \(P \cup (-P) = R\) . Then b is solvable. Let \(b'\) be a solvable subalgebra of g containing b. There exists a subset Q of R such that \(b' = h + g^{Q}\) . Then \(Q \cap (-Q) = \varnothing\) and \(Q \supset P\) , so Q = P and \(b' = b\) .
\end{enumerate}

DEFINITION 1. A subalgebra of g containing h and satisfying the equivalent condition in Prop. 7 is called a Borel subalgebra of (g, h).

A subalgebra \(\mathfrak{b}\) of a splittable algebra \(\mathfrak{g}\) is called a Borel subalgebra of \(\mathfrak{g}\) if there exists a splitting Cartan subalgebra \(\mathfrak{h}'\) of \(\mathfrak{g}\) such that \(\mathfrak{b}\) is a Borel subalgebra of \((\mathfrak{g},\mathfrak{h}')\) .

Let \((\mathfrak{g},\mathfrak{h})\) be a split reductive Lie algebra. Let \(\mathfrak{g}=\mathfrak{c}\times\mathfrak{s}\) with c commutative and s semi-simple. A subalgebra of g of the form \(\mathfrak{c}\times\mathfrak{b}\) , where b is a Borel subalgebra of \((\mathfrak{s},\mathfrak{h}\cap\mathfrak{s})\) , is called a Borel subalgebra of \((\mathfrak{g},\mathfrak{h})\) .

With the notations of Prop. 7, we also say that b is the Borel subalgebra of g defined by h and C (or by h and the basis of R associated to C).

Remark. The map which associates \(\mathrm{R}_{+}(\mathrm{C})\) to a chamber C of R is injective (Chap. VI, §1, no. 7, Cor. 1 of Prop. 20). Consequently, \(\mathrm{C} \mapsto \mathfrak{h} + \mathfrak{g}^{\mathrm{R}_{+}(\mathrm{C})}\) is a bijection from the set of chambers of R to the set of Borel subalgebras of \((\mathfrak{g}, \mathfrak{h})\) . Thus, the number of Borel subalgebras of \((\mathfrak{g}, \mathfrak{h})\) is equal to the order of the Weyl group of R (Chap. VI, §1, no. 5, Th. 2).

PROPOSITION 8. Let \(\mathfrak{b}\) be a subalgebra of \(\mathfrak{g}\) , \(k'\) an extension of \(k\) . Then \(\mathfrak{b} \otimes_k k'\) is a Borel subalgebra of \((\mathfrak{g} \otimes_k k', \mathfrak{h} \otimes_k k')\) if and only if \(\mathfrak{b}\) is a Borel subalgebra of \((\mathfrak{g}, \mathfrak{h})\) .

This is clear from condition (iii) of Prop. 7.

PROPOSITION 9. Let \(\mathfrak{b}\) be the Borel subalgebra of \((\mathfrak{g},\mathfrak{h})\) defined by a chamber C of R. Let \(\mathfrak{n} = \mathfrak{g}^{\mathrm{R}_{+}(\mathrm{C})} = \sum_{\alpha \in \mathrm{R},\alpha >0}\mathfrak{g}^{\alpha}\) . Let \(l = \dim \mathfrak{h}\) .

\begin{enumerate}
\def\labelenumi{(\roman{enumi})}
\item
  If \(h \in \mathfrak{h}\) and \(x \in \mathfrak{n}\) , the characteristic polynomial of \(\mathrm{ad}_{\mathfrak{g}}(h + x)\) is \(\mathrm{T}^l \prod_{\alpha \in \mathbb{R}} (\mathrm{T} - \alpha(h))\) .
\item
  The largest nilpotent ideal of \(\mathfrak{b}\) is equal to \(\mathfrak{n}\) and to \([\mathfrak{b},\mathfrak{b}]\) . This is also the set of elements of \(\mathfrak{b}\) nilpotent in \(\mathfrak{g}\) .
\item
  Let \(\mathrm{B}\) be the basis of \(\mathbb{R}\) associated to \(\mathbb{C}\) . For all \(\alpha \in \mathbb{B}\) , let \(X_{\alpha}\) be a non-zero element of \(\mathfrak{g}^{\alpha}\) . Then \((X_{\alpha})_{\alpha \in \mathbb{B}}\) generates the Lie algebra \(\mathfrak{n}\) . We have \([\mathfrak{n}, \mathfrak{n}] = \sum_{\alpha \in \mathbb{R}, \alpha > 0, \alpha \notin \mathbb{B}} \mathfrak{g}^{\alpha}\) .
\end{enumerate}

There exists a total order on \(\mathfrak{h}_{\mathbf{Q}}^{*}\) compatible with its vector space structure and such that the elements of \(\mathrm{R}_{+}(\mathrm{C})\) are \(>0\) (Chap. VI, §1, no. 7). Let \(h, x\) be as in (i) and \(y \in \mathfrak{g}^{\alpha}\) . Then \([h + x, y] = \alpha(h)y + z\) where \(z \in \sum_{\beta > \alpha} \mathfrak{g}^{\beta}\) . Then, with respect to a suitable basis of \(\mathfrak{g}\) , the matrix of \(\operatorname{ad}_{\mathfrak{g}}(h + x)\) has the following properties:

\begin{enumerate}
\def\labelenumi{\arabic{enumi})}
\item
  it is lower triangular;
\item
  the diagonal entries of the matrix are the number 0 (l times) and the \(\alpha(h)\) for \(\alpha \in \mathbb{R}\) .
\end{enumerate}

This proves (i). It also shows that the characteristic polynomial of \(\operatorname{ad}_{\mathfrak{b}}(h+x)\) is \(\mathrm{T}^{l}\prod_{\alpha\in\mathrm{R}_{+}(\mathrm{C})}(\mathrm{T}-\alpha(h))\) . It follows from the preceding that the set of elements of b nilpotent in g, as well as the largest nilpotent ideal of b, are equal to n.~We have \(n=[b,b]\) by Prop. 5 (i). Finally, assertion (iii) follows from §2, Prop. 4 (ii) and Chap. VI, §1, no. 6, Prop. 19.

COROLLARY. Let \(\mathfrak{b}\) be a Borel subalgebra of \(\mathfrak{g}\) .

\begin{enumerate}
\def\labelenumi{(\roman{enumi})}
\item
  Every Cartan subalgebra of b is a splitting Cartan subalgebra of g.
\item
  If \(\mathfrak{h}_1, \mathfrak{h}_2\) are Cartan subalgebras of \(\mathfrak{b}\) , there exists \(x \in [\mathfrak{b}, \mathfrak{b}]\) such that \(e^{\mathrm{ad}_{\mathfrak{g}} x} \mathfrak{h}_1 = \mathfrak{h}_2\) .
\end{enumerate}

Assertion (i) follows from Prop. 9 (i) and Chap. VII, §3, no. 3, Prop. 3. Assertion (ii) follows from Prop. 9 (ii) and Chap. VII, §3, no. 4, Th. 3.

PROPOSITION 10. Let \(\mathfrak{b},\mathfrak{b}'\) be Borel subalgebras of \(\mathfrak{g}\) . There exists a splitting Cartan subalgebra of \(\mathfrak{g}\) contained in \(\mathfrak{b} \cap \mathfrak{b}'\) .

Let \(\mathfrak{h}\) be a Cartan subalgebra of \(\mathfrak{b}\) , \(\mathfrak{n} = [\mathfrak{b},\mathfrak{b}],\mathfrak{n}' = [\mathfrak{b}',\mathfrak{b}']\) , \(\mathfrak{p} = \mathfrak{b}\cap \mathfrak{b}'\) , and \(\mathfrak{s}\) a vector subspace of \(\mathfrak{g}\) complementary to \(\mathfrak{b} + \mathfrak{b}'\) . Denote by \(\mathfrak{s}^{\perp},\mathfrak{b}^{\perp},\mathfrak{b}'^{\perp}\) the orthogonal complements of \(\mathfrak{s},\mathfrak{b},\mathfrak{b}^{\prime}\) with respect to the Killing form of \(\mathfrak{g}\) . Put \(l = \dim \mathfrak{h}, n = \dim \mathfrak{n}, p = \dim \mathfrak{p}\) . Then \(\dim \mathfrak{b} = \dim \mathfrak{b}^{\prime} = l + n\) ,

\[
\dim \mathfrak {s} ^ {\perp} = \dim (\mathfrak {b} + \mathfrak {b} ^ {\prime}) = 2 (l + n) - p,
\]

and so

\[
\begin{array}{c} \dim (\mathfrak {s} ^ {\perp} \cap \mathfrak {p}) \geq \dim \mathfrak {s} ^ {\perp} + \dim \mathfrak {p} - \dim \mathfrak {g} \\ = 2 (l + n) - p + p - (l + 2 n) = l. \end{array}\tag{3}
\]

By Prop. 1 of §2, no. 2, \(\mathfrak{n} \subset \mathfrak{b}^{\perp}\) , \(\mathfrak{n}' \subset \mathfrak{b}'^{\perp}\) . The elements of \(\mathfrak{p} \cap \mathfrak{n}\) are nilpotent in \(\mathfrak{g}\) (Prop. 9 (ii)), and belong to \(\mathfrak{b}'\) , and hence to \(\mathfrak{n}'\) (Prop. 9 (ii)). Consequently, \(\mathfrak{p} \cap \mathfrak{n} \subset \mathfrak{n} \cap \mathfrak{n}' \subset \mathfrak{b}^{\perp} \cap \mathfrak{b}'^{\perp}\) , so \(\mathfrak{s}^{\perp} \cap \mathfrak{p} \cap \mathfrak{n} = 0\) . In view of (3), we see that \(\mathfrak{s}^{\perp} \cap \mathfrak{p}\) is a complement of \(\mathfrak{n}\) in \(\mathfrak{b}\) . Let \(z\) be an element of \(\mathfrak{h}\) regular in \(\mathfrak{g}\) ; there exists \(y \in \mathfrak{n}\) such that \(y + z \in \mathfrak{s}^{\perp} \cap \mathfrak{p}\) ; by Prop. 9 (i), \(\operatorname{ad}_{\mathfrak{g}}(y + z)\) has the same characteristic polynomial as \(\operatorname{ad}_{\mathfrak{g}}z\) , so \(x = y + z\) is regular in \(\mathfrak{g}\) and a fortiori in \(\mathfrak{b}\) and \(\mathfrak{b}'\) (Chap. VII, §2, no. 2, Prop. 9). Since \(\mathfrak{g}, \mathfrak{b}, \mathfrak{b}'\) have the same rank, \(\mathfrak{b}^{0}(x) = \mathfrak{g}^{0}(x) = \mathfrak{b}'^{0}(x)\) is simultaneously a Cartan subalgebra of \(\mathfrak{b}\) , of \(\mathfrak{g}\) and of \(\mathfrak{b}'\) (Chap. VII, §3, no. 3, Th. 2). Finally, this Cartan subalgebra of \(\mathfrak{g}\) is splitting by the Cor. of Prop. 9.

COROLLARY. The group \(\operatorname{Aut}_{e}(\mathfrak{g})\) operates transitively on the set of pairs \((\mathfrak{t},\mathfrak{b})\) where t is a splitting Cartan subalgebra of g and b is a Borel subalgebra of \((\mathfrak{g},\mathfrak{t})\) .

Let \((\mathfrak{t}_1, \mathfrak{b}_1)\) and \((\mathfrak{t}_2, \mathfrak{b}_2)\) be two such pairs. There exists a splitting Cartan subalgebra \(\mathfrak{t}\) of \(\mathfrak{g}\) contained in \(\mathfrak{b}_1 \cap \mathfrak{b}_2\) (Prop. 10). By the Cor. of Prop. 9, we are reduced to the case in which \(\mathfrak{t}_1 = \mathfrak{t}_2 = \mathfrak{t}\) . Let \(S\) be the root system of \((\mathfrak{g}, \mathfrak{t})\) . There exists bases \(B_1, B_2\) of \(S\) such that \(\mathfrak{b}_i\) is associated to \(B_i\) ( \(i = 1, 2\) ), and there exists \(s \in W(S)\) which transforms \(B_1\) into \(B_2\) . Finally, there exists \(a \in \operatorname{Aut}_e(\mathfrak{g})\) such that \(a|t = s\) ( \(\S 2\) , no. 2, Cor. of Th. 2). Then \(a(t) = t\) and \(a(\mathfrak{b}_1) = \mathfrak{b}_2\) .

